\chapter{绪论}
\label{chap:intro}

\section{研究背景与意义}
\label{sec:background}

近年来，高校智慧校园建设持续推进，教务管理、校园生活、实验室使用和办事流程等事务逐步迁移到线上。然而，与学生日常使用密切相关的信息并没有因此变得集中：培养方案、学籍管理规定、考试安排、宿舍报修、校园网认证、实验室预约等内容，分别由教务、学工、后勤、图书馆和二级学院等不同部门发布，散落在规章制度、通知公告、办事指南和常见问答之中。文件数量多、条款分散、更新节奏不一致，学生往往只能逐页翻阅，或者向辅导员、班委和同学询问。这种依赖人工转述的信息获取方式不仅效率低，而且不同渠道给出的说法有时并不一致，容易造成误解。

通用大语言模型为校园信息获取提供了新的可能。以 Transformer 为基础的大模型能够理解自然语言问题，并以接近日常交流的方式组织答案，一定程度上降低了用户表达问题和理解答案的门槛。但通用模型的知识主要来自公开训练语料，并不掌握学校最新的校本文件；直接让它回答“补办学生证需要哪些材料”“某门课程重修如何申请”等问题，模型可能给出看似合理却没有依据的内容，即通常所说的“幻觉”。更关键的是，通用模型一般不会说明答案出自哪一份文件、哪一条规定，用户难以核实，学校也难以对回答质量进行追责。检索增强生成（Retrieval-Augmented Generation，RAG）的出现为这一问题提供了可行路径：系统先从外部知识库中检索与问题相关的证据，再把证据和问题一起交给大模型，让模型在给定材料的范围内组织回答\upcite{lewis2020}。在校园场景中引入 RAG，可以把“查文件”和“问模型”结合起来，使回答既有自然语言的易读性，又能追溯到具体资料\upcite{baijing2026}。

本课题的理论意义在于，把中文校园文档的分块、混合检索、重排序、受约束生成和引用溯源串成一条完整流程，并通过对照实验观察检索增强前后回答准确性和可核查性的变化，为 RAG 在校园垂直场景中的落地提供可复用的设计参考。实践意义则体现在三个方面：一是降低学生获取校园信息的成本，减少在多个文件之间反复查找的时间；二是减轻辅导员和管理人员的重复答疑负担；三是通过问答历史、高频问题和满意度反馈，为学校改进制度宣传和服务流程提供数据依据。与训练或微调专用模型相比，基于成熟大模型接口和开源检索组件的方案实施成本较低，便于在本科毕业设计的周期内完成，也便于后续推广到其他高校。

\section{国内外研究现状}
\label{sec:related}

\subsection{检索增强生成的研究进展}

检索增强生成的核心思想，是把参数化的大模型与非参数化的外部知识库结合起来。Lewis 等提出的 RAG 框架由检索器和生成器两部分组成：检索器从维基百科的稠密向量索引中取出相关段落，生成器以问题与段落为条件产生答案\upcite{lewis2020}。这一工作说明，回答所需的证据不必全部固化在模型参数中，可以在推理阶段从外部资料现取。此后，围绕 RAG 的优化逐渐形成一条较完整的技术路线。袁乐等将 RAG 的流程拆分为分块、向量化、建索引、查询处理、块选择和提示构建等环节，并从预检索增强、检索器增强、检索策略增强、索引增强、检索后增强和大模型增强六个方面归纳了优化方法\upcite{yuan2026}。张犬俊等则从软件工程的角度梳理了 RAG 在需求分析、代码生成、测试和维护等任务中的应用，指出检索质量与上下文组织方式对最终效果具有决定性影响\upcite{zhang2026}。

在生成控制方面，Asai 等提出的 Self-RAG 让模型按需发起检索，并通过反思标记评价检索段落与自身生成内容的相关性和事实性\upcite{asai2024}。该方法在多个知识密集型任务上取得了优于常规 RAG 的表现，但需要额外的训练和多次模型调用，成本较高。对于本科课题而言，更现实的做法是借鉴其“先判断证据是否充分”的思想，通过提示词约束、证据阈值和拒答策略实现轻量化的可控生成。总体来看，RAG 已经形成较为清晰的技术框架，但面向校园这类资料格式多样、更新频繁、权限要求明确的场景，仍需要针对性的工程化设计。

\subsection{检索与重排序技术}

检索环节决定了送入生成模型的证据质量。早期的问答系统多采用关键词匹配或 BM25 等稀疏检索方法，这类方法对专有名词、文件编号和条款序号比较敏感，但难以处理同义表达和语义泛化。Karpukhin 等提出的密集段落检索（Dense Passage Retrieval，DPR）使用双塔结构，将问题和段落分别编码为稠密向量，再通过向量相似度召回候选段落\upcite{karpukhin2020}。与 BM25 相比，DPR 不再依赖词语是否重合，能够在一定程度上理解语义相近但措辞不同的问法。中文场景下的向量化则可以直接借助已有资源。Xiao 等构建的 C-Pack 包含中文评测基准、大规模训练语料和 BGE 系列预训练模型，为中文语义表示提供了成套工具\upcite{xiao2024}。本课题选用 BGE 系列模型完成知识片段和问题的向量化，避免从零训练嵌入模型。

单一检索通道各有短板。郜洪奎等指出，单纯依赖向量表示难以覆盖检索任务的复杂性，而且缺少有效的排序机制，容易把无关信息带入上下文；他们先用多路径召回扩大候选范围，再用双层网络对候选进行多轮排序\upcite{gao2025}。张健等发现，按段落切分比按固定长度切分更有利于保持语义完整，检索知识放在问题之前时采用逆序重排序可以改善回答准确率\upcite{zhangjian2025}。孙飒等面向企业知识库提出结构感知的自适应分块、由摘要与细粒度内容构成的层次化父子块索引，以及双通道混合检索和交叉重排方案，以解决固定窗口切分把完整操作步骤拆散、扁平化索引丢失页面层级的问题\upcite{sun2026}。校园文件同样具有标题层级和条款结构，这些工作为本课题的分块策略和混合检索设计提供了直接参考。

\subsection{垂直领域问答系统的应用}

在校园和垂直领域，RAG 已经被用于多项知识密集型问答任务。白晶等把高校知识问答系统自下而上划分为数据层、检索层、模型层和应用层，并建立了覆盖性能指标、用户体验和知识质量的三维评价体系\upcite{baijing2026}。宋保仁和冷静以四川高校图书馆特色数据库为例，分析了 RAG 在私有知识库建设中的应用与边界，认为该技术提升了特色资源的检索效率和知识组织水平，同时提醒要配套可审查的生成路径和协同治理机制\upcite{song2026}。这一结论对本课题具有警示意义：校园资料的范围、更新责任和答案可核查性，必须在系统设计阶段就明确下来。

秦晓瑞等面向结构复杂、语义严谨的航空标准文档，设计了结构路径感知的知识库，用章节结构和标题链条支持语境追溯，并在检索阶段融合关键词与语义信号，由大模型判断是否需要拆解子问题、发起多轮检索\upcite{qin2026}。甘克勤等把对话历史管理单独做成模块，通过主题建模和序列编码分析会话的对象、主题和结构，使系统在多轮问答中能够理解指代和上下文\upcite{gan2026}。罗云等提出的 SV-RAG 在切分时保留层次结构元数据，检索时采用双轨混合检索，生成时要求模型先引用核心原文再展开说明，以增强回答的忠实性\upcite{luo2026}。这种“先引用、后阐述”的顺序，对本课题的引用溯源设计很有启发。

此外，董永涛等设计了舰艇装备故障智能问答系统，彭琳讨论了企业级知识库的构建与优化，陶金涛等将 RAG 用于新疆地质矿产问答，王祥昊等将其用于铁路应急处置策略生成，窦嘉铭等则把知识图谱与 RAG 结合，用于配网检修辅助决策\upcite{dong2025,peng2026,tao2026,wang2026,dou2026}。这些工作说明，RAG 在专业资料密集、答案准确性要求高的场景中具有较好的适用性；但校园场景的特殊性在于资料以中文行政文本为主、更新频繁、用户角色多样，现有研究对这类资料的持续建库和权限隔离讨论得还不够充分。

\subsection{评价、安全与工程实现}

如何评价 RAG 系统的效果，是另一条重要线索。Es 等提出的 RAGAS 不依赖人工标注的参考答案，围绕忠实性、答案相关性和上下文相关性三个维度自动打分，分别对应答案能否由上下文推出、是否回答了问题、检索内容是否集中于问题\upcite{es2024}。这套指标为本课题的效果测试提供了参考。不过，校园问答还需要关注引用是否真实、证据不足时能否拒答，因此仅靠通用指标并不够，需要结合人工标注的校园问题测试集进行补充评价。

在系统安全和工程实现方面，卢万有给出了基于 JWT 的无状态认证与 RBAC 角色权限控制方案，身份信息放在令牌中在前后端之间传递，后端通过拦截器完成校验\upcite{lu2025}。廖明焱针对企业 RAG 的权限控制缺失和敏感信息泄露问题，设计了多层级权限控制的安全检索架构与智能脱敏方法\upcite{liao2026}。魏巍则从云资源管理的角度讨论了 RBAC 模型的权限设计\upcite{wei2025}。这些工作说明，知识库系统的权限控制不能只做在页面入口，还要落实到检索和证据返回环节。张宁以 LangChain 为框架，接入本地部署的 DeepSeek 模型，用 Chroma 存放知识向量，前后端采用 Vue 3 和 FastAPI，其技术选型与本课题较为接近\upcite{zhangning2026}。吴海周等把系统分为索引、检索、生成三个模块分别优化，采用语义向量构建、动态 Top-$k$ 检索和上下文融合生成等方法，说明检索条数、上下文长度等参数需要结合具体数据调整\upcite{wu2026}。LangChain 官方文档则为文档加载、切分、向量存储和检索链的组织提供了工程参考\upcite{langchain2026}。

\subsection{现有研究的不足}

综合上述工作可以看出，RAG 在知识密集型问答中的有效性已经得到较多验证，但放到校园场景中仍有三个问题没有很好解决。第一，校园资料来自不同部门，格式不统一、更新频繁，现有研究对这类资料如何持续建库、如何保留标题层级和条款结构讲得不够细。第二，不少系统更关注回答读起来是否顺畅，对引用是否真实、证据不足时能否拒答、不同角色的资料是否需要隔离考虑得较少。第三，效果评价大多沿用通用指标，缺少针对本校常见问题和引用正确性的测试集。本课题围绕这三点展开：以学校公开资料建立可控的校园知识库，用关键词检索和向量检索双路召回后再重排，生成时限定只能使用检索到的证据并标出引用编号，最后从检索命中、答案正确、引用正确、拒答准确和响应时延等方面设计验证方案。

\section{研究内容与方法}
\label{sec:content}

本课题的研究内容可以概括为四个方面。第一，校园知识库构建。研究 PDF、Word 和纯文本等常见格式的解析方法，设计保留标题层级、页码、分类和更新时间的结构感知分块策略，并完成文本向量化和索引构建。第二，混合检索与重排序。在 BM25 稀疏检索和 BGE 稠密向量检索两条通道上分别召回候选，使用倒数排序融合（Reciprocal Rank Fusion，RRF）合并结果，再通过交叉编码器对顶部候选进行重排。第三，带引用的受约束生成。设计包含问题、证据和输出约束的提示词模板，要求模型只依据证据回答、为关键结论标注引用编号，并在证据不足时拒答；同时处理多轮对话中的指代和追问。第四，系统实现与效果评价。采用 Vue 3、FastAPI、MySQL 和向量数据库完成 Web 系统，设计功能测试、性能测试和对照实验，比较纯大模型问答与不同检索策略的差异。

研究方法上，本文采用文献分析、需求调研、原型开发和对照实验相结合的方式。文献分析用于梳理 RAG、稠密检索、混合检索、重排序和评价指标的研究基础；需求调研结合任务书、开题报告和校园常见问题，明确学生与管理员的使用场景；原型开发按照软件工程的流程完成需求分析、总体设计、详细设计、编码实现和联调；对照实验则在同一测试集上比较不同方案，尽量保证变量可控、结果可复现。

\section{论文组织结构}
\label{sec:organization}

本文共分为六章，另设结论、参考文献和致谢。第 1 章介绍研究背景、国内外研究现状、研究内容与方法，说明课题的来源和意义。第 2 章介绍大语言模型、检索增强生成、文本向量化、混合检索与重排序、多轮对话和 Web 权限控制等关键技术，为后续设计提供理论依据。第 3 章从可行性、用户角色、功能需求、非功能需求和业务流程等方面完成系统需求分析。第 4 章给出系统总体架构、知识库与数据处理、检索增强问答流程、数据库与接口、权限与安全等设计。第 5 章说明开发环境、项目结构以及用户权限、文档管理、混合检索、受约束生成和管理统计等模块的实现。第 6 章给出测试环境、功能与性能测试、问答测试集与评价指标、对照实验和结果分析方法。结论部分总结全文工作，指出不足并展望后续改进方向。
